% META: {"id":"1SPE-EXPO-EX-028","chapitre":"1SPE-EXPONENTIELLE","type_objet":"exercice","parcours":2,"competences":["raisonner"],"duree_min":10,"mode_creation":"ex_nihilo","sources_inspiration":[],"fichier_tex":"chapitres/1SPE-EXPONENTIELLE/exercices/1SPE-EXPO-EX-028.tex","corrige_tex":"chapitres/1SPE-EXPONENTIELLE/corriges/1SPE-EXPO-CO-028.tex","status":"approved","extension_codes":["X1"],"programme_alignment":"OPTIONAL_EXTENSION","extension_label":"Approfondissement — Vers la Terminale","origin":"generated","status_history":["generated","audited","accepted","approved"],"release_acceptance":"RELEASE_OWNER_BATCH_ACCEPTANCE","acceptance_closure_digest":"sha256:9b3ccf9a81c5520fbb7e03b7b2d3e7bf2057a02834b6d908bf3be5f49f3b553f","acceptance_receipt_digest":"sha256:4fad86f0282e0fa4e0419cca1ad6379ae7af5a280c04a4a90880f90a1c044242"}
\begin{center}\textbf{Approfondissement — Vers la Terminale}\end{center}
% BEGIN-VERIFY
% from sympy import *
% x = symbols('x')
% g = exp(x) - x - 2
% assert g.subs(x, 0) == -1
% assert g.subs(x, 2) == exp(2) - 4
% assert g.subs(x, 0) < 0
% assert g.subs(x, 2) > 0
% END-VERIFY
\begin{exercice}{1SPE-EXPO-EX-028}{2}{10}
On pose $g(x) = \mathrm{e}^x - x - 2$ pour tout réel $x$.
\begin{enumerate}
  \item Calculer $g(0)$ et $g(2)$.
  \item Montrer que $g$ est continue sur $[0\,;\,2]$.
  \item En déduire que l'équation $\mathrm{e}^x = x + 2$ admet au moins une solution dans l'intervalle $[0\,;\,2]$.
\end{enumerate}
\end{exercice}
